\chapter{Objetivos}
\label{cap:objetivos}
\section{Objetivo general}
Diseñar, implementar y evaluar una plataforma de entrenamiento auditivo musical basada
en juegos serios que integre el análisis de respuestas vocales, la retroalimentación al
estudiante y el ajuste de la dificultad de los ejercicios según su desempeño.

Este objetivo recupera la formulación del protocolo presentado el 28 de septiembre de
2026. Su evaluación se centra en la precisión de la detección de tono, la latencia de la
retroalimentación y el comportamiento del ajuste de dificultad en dispositivos móviles.
\section{Objetivos específicos}
Los siguientes objetivos específicos descomponen el objetivo general y el producto
esperado del protocolo, en correspondencia con las tareas del plan adjunto. Constituyen
una formulación de trabajo revisada con el autor, pendiente de revisión con el director.

\begin{description}
  \item[OE1.] Diseñar e implementar un módulo de captura y análisis de audio monofónico
  en el dispositivo que estime el tono de las respuestas vocales y permita presentar
  retroalimentación respecto de un objetivo musical.
  \item[OE2.] Diseñar e implementar un mecanismo configurable de ajuste de dificultad
  que utilice indicadores de desempeño y pueda integrarse en distintas actividades.
  \item[OE3.] Integrar el análisis vocal, la retroalimentación y el ajuste de dificultad
  en un prototipo Android con un ejercicio de reproducción vocal de intervalos y otro
  minijuego de entrenamiento auditivo, conforme al mínimo establecido en el plan de
  trabajo. Explicitar la habilidad musical que practica cada actividad y contemplar
  variaciones pertinentes para reducir la dependencia de pistas propias del juego.
  \item[OE4.] Incorporar seguimiento del progreso y registro exportable de datos de
  desempeño y funcionamiento que permitan analizar las respuestas y las decisiones
  de adaptación.
  \item[OE5.] Evaluar la precisión de la estimación de tono, la latencia de la
  retroalimentación y el comportamiento del ajuste de dificultad mediante pruebas
  con referencias, condiciones de ejecución y criterios explícitos.
  \item[OE6.] Examinar el uso del prototipo mediante sesiones exploratorias inicialmente
  con estudiantes de música mayores de edad con conocimientos básicos, y documentar
  las dificultades observadas y las mejoras derivadas, delimitando las conclusiones
  a las condiciones y participantes evaluados.
  \item[OE7.] Publicar una versión del prototipo gratuita en todas sus funciones,
  sin suscripción y utilizable sin conexión, junto con su código fuente y su
  documentación bajo una licencia de código abierto compatible con los recursos
  y las dependencias utilizados.
\end{description}
